% \vspace{-2mm}
\section{Conclusion and Limitation}
% \vspace{-2mm}
In this work, we conducted a systematic investigation to decouple the effects of modeling paradigms (AR vs. MDM) from their commonly associated architectural choices (decoder-only vs. encoder-only). By implementing MDMs (as AO-AR) within a decoder-only framework, we facilitated a more equitable comparison and explored architectural influences within the MDM paradigm itself.
Our comparison of autoregressive (AR) and masked diffusion models (MDMs) within a decoder-only framework revealed that MDM's uniform order-agnosticism may be suboptimal for language, given its inherent left-to-right structure. This suggests future MDM research could benefit from exploring non-uniform order distributions, balancing modeling power with data alignment and efficiency. Furthermore, contrasting encoder-only and decoder-only MDMs highlighted decoders' significantly lower generation complexity (e.g., linear vs. quadratic), though encoders offer unique advantages like bidirectional attention and context order invariance. These insights underscore the need to consider architectural impacts beyond formulation when comparing models, and affirm the strong potential of decoder-based MDMs as an efficient direction for future exploration.

Our experiments were conducted on models of up to medium size (e.g., ~350M parameters). Whether these observations generalize to significantly larger computational scales remains an open question. Furthermore, this work focused on language; the applicability of our findings to other discrete data modalities is uncertain, especially as many such modalities may not possess the strong left-to-right sequential structure inherent in natural language.
